% ============ 第2章 §2.1 Hardy空间的定义(原书 pp.124-126) ============
\section{Hardy 空间的定义}\label{sec:2.1}
\index{Hardy 空间!(单位圆板上的)}


$\mathbb{C}$ 为复平面,$D = \{z \in \mathbb{C};\ |z| < 1\}$,$\partial D = \{z \in \mathbb{C};\ |z| = 1\}$,且 $\overline{D} = \{z \in \mathbb{C};\ |z| \leq 1\}$.$h(D)$ 表示 $D$ 上实值调和函数的全体,$H(D)$ 表示 $D$ 上正则函数的全体.对 $f \in H(D)$,$\operatorname{Re}(f)$ 表示 $f$ 的实部,$\operatorname{Im}(f)$ 表示 $f$ 的虚部.令 $\operatorname{Re} H(D) = \{\operatorname{Re}(f);\ f \in H(D)\}$,则因 $D$ 是单连通的,$\operatorname{Re} H(D) = h(D)$ 成立.又 $\operatorname{Im} H(D) = \{\operatorname{Im}(f);\ f \in H(D)\} = h(D)$.为研究属于 $h(D)$ 与 $H(D)$ 的函数,可以考察多种度量函数在 $D$ 的边界附近绝对值的增长度的量.特别地,当 $F$ 在 $D$ 上连续时,
\begin{equation*}
M_{p}(r, F) = \left\{\frac{1}{2\pi}\int_{-\pi}^{\pi}|F(re^{i\theta})|^{p}\dd\theta\right\}^{1/p} \quad (0 < p < \infty)
\end{equation*}
与
\begin{equation*}
M_{\infty}(r, F) = \max_{-\pi \leq \theta < \pi}|F(re^{i\theta})|
\end{equation*}
是应用上重要的度量 $F$ 的增长度的量.当然,当 $F$ 在 $\overline{D} = \{z \in \mathbb{C};\ |z| \leq 1\}$ 上连续时,对任意 $p$,$\sup_{0 \leq r < 1}M_{p}(r, F) < \infty$.这里对实值调和函数与正则函数定义 Hardy 空间 $h^{p}$ 与 $H^{p}$.对 $0 < p \leq \infty$,
\begin{equation*}
h^{p} = \left\{u \in h(D);\ \sup_{0 \leq r < 1}M_{p}(r, u) < \infty\right\}
\end{equation*}
且
\begin{equation*}
H^{p} = \left\{f \in H(D);\ \sup_{0 \leq r < 1}M_{p}(r, f) < \infty\right\}
\end{equation*}
由定义,$h^{\infty} \subset h^{p} \subset h^{q}$($p > q$)且 $H^{\infty} \subset H^{p} \subset H^{q}$($p > q$)成立.$h^{\infty}$ 与 $H^{\infty}$ 显然是线性空间.当 $0 < p < \infty$ 时,
\begin{equation*}
a^{p} \leq (a + b)^{p} \leq 2^{p}(a^{p} + b^{p}) \quad (a \geq 0, b \geq 0)
\end{equation*}
故 $h^{p}$ 与 $H^{p}$ 是线性空间.虽然 $\operatorname{Re} H(D) = h(D)$,但 $\operatorname{Re} H^{p} = h^{p}$ 是否成立,对一般的 $p$ 是难以证明的.多项式 $\sum_{n=0}^{\ell}a_{n}z^{n}$ 属于 $H^{\infty}$,$\sum_{n=0}^{\ell}a_{n}z^{n} + \sum_{n=0}^{\ell}\overline{a_{n}z^{n}}$ 属于 $h^{\infty}$.当 $0 < p < 1$ 时 $H^{p}$ 作为空间是很大的,$(1-z)^{-1}$ 不属于 $H^{1}$,但可以证明 $(1-z)^{-1} \in \bigcap_{0 < p < 1}H^{p}$.另外,$\operatorname{Re}(1-z)^{-1}$ 与 $\operatorname{Re}(1+z)(1-z)^{-1}$ 属于 $h^{1}$ 这一点将在后面证明.

下一个引理中的不等式在 Hardy 空间 $H^{p}$ 的研究中将在后面用到.

\begin{lemma}\label{lem:2.1.1}
设 $0 < p \leq \infty$.若 $f \in H(D)$ 且 $0 < r < r' < 1$,则 $M_{p}(r, f) \leq M_{p}(r', f)$.
\end{lemma}

\begin{proof}
设 $p \neq \infty$.取 $u(z)$ 为在 $|z| < r'$ 上调和,在 $|z| \leq r'$ 上连续,且在 $|z| = r'$ 上等于 $|f(z)|^{p}$ 的函数.由于 $|f|^{p}$ 是次调和函数,在 $|z| \leq r'$ 上 $|f(z)|^{p} \leq u(z)$.因此,利用调和函数的 Gauss 平均值定理,
\begin{align*}
M_{p}(r, f) & \leq \frac{1}{2\pi}\int_{-\pi}^{\pi}u(re^{i\theta})\dd\theta = u(0) \\
& = \frac{1}{2\pi}\int_{-\pi}^{\pi}u(r'e^{i\theta})\dd\theta = M_{p}(r', f).
\end{align*}
当 $p = \infty$ 时,由正则函数的最大值原理是显然的.
\end{proof}
